%%%%%%%%%%%%%%%%%%%%%%%%%%%%%%%%%%%%%%%%%%%%%%
%% 03_shaping_the_lookback.tex --- 正文层 (BODY LAYER)
%%
%% 这里只有内容。字体、边距、颜色、定理环境长什么样，全在
%% ../../preamble.tex 里；改版式不用碰这个文件。
%%
%% 这是正文片段，没有 \documentclass 也没有 document 环境：
%% ../../build.sh 单独编译它，all_chapters.tex 把它 \input 进合订本。
%%%%%%%%%%%%%%%%%%%%%%%%%%%%%%%%%%%%%%%%%%%%%%

\chapterhead{03}{Shaping the Lookback}{03_shaping_the_lookback}

\chapref{2 \midot{} Testing a Signal}{02_testing_a_signal} left two blanks inside its own
definition. Writing momentum as $\operatorname{Avg}(r_{s,t-i})$, $i = 1 \ldots N$ fixes neither the
length of the window nor the weight each return carries inside it---a flat average gives every
return in the window an equal vote regardless of age, and commits to one window that cannot be both
stable and timely.

An \textbf{exponentially weighted moving average (EWMA)} answers the first blank: replace the flat
weights with a decay, so the newest return counts for the most and the average never quite forgets
(\secref{sec:ewma}). \textbf{Moving average convergence/divergence (MACD)} answers the second blank
on top of that fix---instead of one EWMA it carries two, a fast one and a slow one, and trades the
gap between them (\secref{sec:fast-slow}, \secref{sec:macd}). MACD is therefore not a third
technique standing beside EWMA; it is EWMA's own decay run twice, at two speeds, and read as a
difference.

Both pieces run into the same wall no matter how they are tuned, and \secref{sec:vol-clustering} is
that wall: \textbf{a failure mode that survives every setting of either}, because it belongs to the
market rather than to the parameters. The repair \secref{sec:vol-clustering} can offer is local; the
repair it cannot offer is \chapref{4 \midot{} Volatility Regimes}{04_volatility_regimes}.


\section{EWMA instead of a simple moving average}\label{sec:ewma}

\subsection{Why the newest return should weigh most}\label{sec:newest-weighs-most}

A momentum signal is a weighted sum of past returns, and a simple moving average is one particular
choice of those weights: make every one of them equal. Stated that way it is a claim about
information---that a return from sixty days ago tells you exactly as much about where the asset
stands today as yesterday's does.

It does not. \textbf{The newest return is the most recent evidence about the state the asset is in
now}, so it should carry the most weight, and a return from three months ago the least. An
exponentially weighted moving average is that preference written down.

\fig{ewma-weights}{Weight carried by each past return against its lag in trading days, both curves
normalised to the same total. The 21-day simple moving average is a flat box---identical weight out
to twenty-one days, then nothing. The EWMA starts highest at lag zero, decays smoothly, and stays
above zero past forty-five days. Illustrative.}{fig:ewma-weights}

Two things change at once. Relative to the box the newest lags are \textbf{over-weighted} and the
old ones \textbf{under-weighted}---and the box's hard edge at twenty-one days is gone, so an EWMA
thins out without ever quite forgetting.

\subsection{The recursion, and the weights it produces}\label{sec:ewma-recursion}

\begin{definition}[EWMA]\label{def:ewma}
Carry one running number and update it as each observation arrives:
\[
  \mathrm{EWMA}_t = (1 - \lambda)\, a_t + \lambda\, \mathrm{EWMA}_{t-1}
\]
where $a_t$ is the newest observation and $\lambda \in (0, 1)$ the decay---how much of the old
average survives each step.
\end{definition}

\begin{example}\label{ex:ewma-half}
Take $\lambda = 1/2$ and three observations $a_1, a_2, a_3$, oldest first. Two updates run the whole
thing:
\begin{align*}
  \text{after } a_2 &: \quad \frac{a_1 + a_2}{2} \\
  \text{after } a_3 &: \quad \frac{1}{2} \cdot \frac{a_1 + a_2}{2} + \frac{a_3}{2}
                       \;=\; \frac{a_3}{2} + \frac{a_2}{4} + \frac{a_1}{4}
\end{align*}

\begin{center}
\begin{tabular}{|c|l|c|}
\hline
\textbf{Observation} & \textbf{Age} & \textbf{Weight it ends up with} \\ \hline
$a_3$ & newest        & $\mathbf{1/2}$ \\ \hline
$a_2$ & one step back & $1/4$ \\ \hline
$a_1$ & oldest        & $1/4$ \\ \hline
\end{tabular}
\end{center}

The newest observation carries half the total on its own, and each step back halves what is
left---which is what \emph{exponential} names. The oldest two tie only because the recursion had to
start somewhere; run it long enough and that seam washes out.
\end{example}

\subsection{Choosing how fast it decays}\label{sec:half-life}

Tune the \textbf{half-life} $H$---the lag at which a return's weight has decayed to half the weight
given to the newest one. Candidates are fractions of the lookback ($1/2$, $1/3$, $1/5$, $1/8$),
and---like every free parameter in this chapter---they are grid-searched rather than chosen;
\secref{sec:pairing} grid-searches the analogous ratio for the fast/slow pairing
\secref{sec:fast-slow} builds. In pandas the whole signal is
\texttt{returns.ewm(halflife=H).mean()}.

No value of $H$ is standard, and whatever a package ships as its default is an \textbf{empirical
solution}---a number that fit somebody's historical data once. Every parameter in this chapter has
that status, which is \chapref{5 \midot{} Modelling}{05_modelling}'s subject rather than something
to accept on authority.


\section{Combining a fast and a slow horizon}\label{sec:fast-slow}

\secref{sec:ewma} fixed the weight-shape blank with a decay. This section fixes the other one---not
by finding a better single window, but by refusing to settle for just one.

\subsection{Why two windows beat one}\label{sec:two-windows}

One lookback forces a choice between stable and timely. A long window rides a trend without being
shaken out of it but arrives late at both ends; a short one turns on time, and turns on noise too.
Rather than choose, carry both and trade \textbf{the gap between them}.

That gap is a single series---the fast average minus the slow one---and the rule is its sign: long
while the fast average sits above the slow one, out when it drops back under. Two windows go in and
one number comes out, which is what makes the pair an indicator rather than two rules run side by
side.

\fig{fast-times-slow}{Above, one price path with a 20-day and a 40-day moving average over it; the
fast average crosses above the slow one, and later drops back under, and the bar between those two
dates marks the stretch the rule is long. Below, the single series the pair produces---fast minus
slow---against zero, shaded where it is positive. Its two zero crossings fall exactly on the two
dates the averages cross above. Illustrative.}{fig:fast-times-slow}

\textbf{What the crossing means.} The fast average overtaking the slow one is not only an arithmetic
consequence of the shorter window. It is the first evidence that flow has changed direction---that
whoever was selling has stopped, or something large has started buying---while the move is still too
small to register over a quarter. The same reading runs the other way and matters more there: the
fast average rolling back under the slow one is what a large holder beginning to leave looks like
from outside. Waiting for the slow window to turn on its own means selling to them on the way out.

\subsection{Finding the pairing}\label{sec:pairing}

The pairing is found by experiment, not chosen. Put the slow lookback $N_s$ on one axis and the fast
one $N_f$ on the other, run \chapref{2}{02_testing_a_signal}'s bar-plot test in every cell, and read
the grid as a heat map.

\fig{ratio-grid}{A five-by-five grid with a geometric ladder of lookbacks---one day, one week, two
weeks, one month, one quarter---on both axes: fast across, slow up. Every cell on and above the
diagonal is struck out, the fast leg having to be the shorter one. Of the rest, the darkest run in a
diagonal band, and the two annotated cells sit at roughly two to one.
Illustrative.}{fig:ratio-grid}

Two things are true of that grid before any data is collected. \textbf{Half of it is empty by
construction}---$N_f < N_s$ or there is no fast leg, so everything on and above the diagonal is
struck out. And the candidates should be spaced \textbf{geometrically} rather than evenly: a day, a
week, two weeks, a month, a quarter. Even spacing spends most of the grid on pairs that differ by a
rounding error.

\textbf{What a diagonal band means.} The warm cells run diagonally, and that shape is the finding
rather than any single cell in it. On a geometric ladder a diagonal is a constant \textbf{ratio}, so
what carries is $N_s / N_f$ and not either window on its own---which is why the answer is quoted as
a ratio at all. Research lands it near \textbf{2:1}, and MACD's 26/12 is one instance; for
daily-to-weekly holding the band shows up around two weeks against a month. An empirical
regularity, not a law, and \chapref{5 \midot{} Modelling}{05_modelling}'s subject the moment you
start trusting it.

\begin{note}[One hot cell is not a band]
A single dark square with pale neighbours is an artifact of however many parameter combinations were
tried, not an edge. What earns belief is a contiguous warm region, because a real effect cannot
switch off between one week and eight days.
\end{note}

\subsection{Where the combination pays}\label{sec:where-it-pays}

Best in \textbf{commodities}, and the reason is what moves them. An equity price answers to
earnings, to the firm's own performance, to whoever is running it---many small idiosyncratic forces,
none of which trend for long. A commodity answers to global energy supply, shipping, and whether the
large participants are in the market at all, and it is far too large for one fund to push on its
own, so a move that starts tends to be everyone moving together. Equities second. \textbf{Bonds}
weakest: lean on one hard enough and someone takes the other side and presses it back.


\section{MACD, stated precisely}\label{sec:macd}

\secref{sec:ewma} gave a window a decay; \secref{sec:fast-slow} took the difference of a fast
average and a slow one and read its sign. That difference already has a name, and two further series
conventionally travel with it.

\subsection{The three series}\label{sec:three-series}

Written out, MACD is three series built from two exponential moving averages of the \textbf{price}.

\begin{definition}[MACD]\label{def:macd}
For a fast and a slow span $n_f < n_s$, each with its own smoothing constant
$\alpha = 2/(\text{span} + 1)$:
\begin{align*}
  \mathrm{MACD}_t   &= \mathrm{EMA}_{n_f}(P)_t - \mathrm{EMA}_{n_s}(P)_t \\
  \mathrm{signal}_t &= \mathrm{EMA}_9(\mathrm{MACD})_t \\
  \mathrm{hist}_t   &= \mathrm{MACD}_t - \mathrm{signal}_t
\end{align*}
where $P_t$ is the price on date $t$ and $\mathrm{EMA}_n(P)_t$ its exponential moving average over
span $n$. The asset subscript is dropped throughout---MACD is computed one asset at a time.
Conventionally $n_f = 12$, $n_s = 26$, and 9 for the signal line.
\end{definition}

\begin{tabularx}{\linewidth}{|>{\raggedright\arraybackslash}p{0.20\linewidth}|X|>{\raggedright\arraybackslash}p{0.26\linewidth}|}
\hline
\textbf{Series} & \textbf{Reads as} & \textbf{Sign means} \\ \hline
\textbf{MACD line}   & recent average price versus a longer one & trend direction \\ \hline
\textbf{Signal line} & a smoothed MACD                          & the level MACD is crossing \\ \hline
\textbf{Histogram}   & MACD minus its own smoothing             & trend acceleration \\ \hline
\end{tabularx}

\subsection{Why it is momentum, not a separate indicator}\label{sec:macd-is-momentum}

\begin{claim}\label{clm:macd-is-momentum}
MACD is a momentum signal. It is a weighted sum of past returns, differing from a lookback mean only
in the shape of the weights.
\end{claim}

\begin{proof}
Each EMA is a weighted average of past prices whose weights sum to one, so their difference weights
the price $i$ days ago by $c_i$, and those weights sum to \textbf{zero}:
\[
  \mathrm{MACD}_t = \sum_i c_i P_{t-i}, \qquad
  c_i = \alpha_f (1-\alpha_f)^i - \alpha_s (1-\alpha_s)^i, \qquad \sum_i c_i = 0 .
\]
A zero-sum weighting is unchanged when every price shifts by a constant, so subtract $P_t$ from each
term and write $P_t - P_{t-i}$ as a sum of one-period changes
$\Delta_{t-j} = P_{t-j} - P_{t-j-1}$. Collecting the coefficient of the change at lag $j$ gives the
\textbf{kernel} $k_j$:
\[
  \mathrm{MACD}_t = \sum_j k_j \Delta_{t-j}, \qquad
  k_j = \sum_{i \leq j} c_i = (1-\alpha_s)^{j+1} - (1-\alpha_f)^{j+1} .
\]
Since $\alpha_f > \alpha_s$, every $k_j \geq 0$: MACD is a non-negative weighted sum of past price
changes, exactly like a lookback mean.
\end{proof}

\begin{note}[What actually differs]
The kernel. A 21-day momentum weights the last 21 returns equally and everything older at zero;
MACD's weights rise from a small value at lag 0, peak around lag 8, and decay without ever reaching
zero. It therefore discounts \emph{yesterday} relative to last week---deliberately, since the newest
return is the noisiest---and never fully forgets.
\end{note}

\fig{signal-kernels}{Weight given to each past daily return against its lag in trading days, for two
signals normalised to the same total. A 21-day momentum is a flat box: equal weight for the last 21
days and zero beyond. MACD with spans 12 and 26 is a hump that starts low at lag zero, peaks around
lag 8, then decays slowly and never reaches zero within sixty days.
Illustrative.}{fig:signal-kernels}

\begin{note}[A choice, not a theorem]
\secref{sec:newest-weighs-most} argued that the newest return deserves the most weight, and the most
widely used trend indicator in the market does the opposite. Both positions are defensible: the
newest return is the most recent evidence about the state the asset is in, and it is also the one
carrying the most reversal (\chapref{2}{02_testing_a_signal}'s Background). Which wins is an
empirical question settled by the bar plot, not a matter of taste---which is why
\secref{sec:newest-weighs-most} was written as a preference rather than a result.
\end{note}

\subsection{From a value to a rule}\label{sec:value-to-rule}

Three common rules, in increasing order of information kept. All three still have to pass
\chapref{\S 2.3.2}{02_testing_a_signal}'s bar plot before they earn a backtest.

\begin{tabularx}{\linewidth}{|>{\raggedright\arraybackslash}p{0.16\linewidth}|>{\raggedright\arraybackslash}p{0.30\linewidth}|X|}
\hline
\textbf{Rule} & \textbf{Go long when} & \textbf{Costs} \\ \hline
Zero-line    & MACD is above zero
             & a slow trend filter, late \\ \hline
Crossover    & the histogram is above zero
             & earlier, noisier---the churn is \secref{sec:vol-clustering}'s subject \\ \hline
Proportional & always, sized by the standardized value
             & none of the magnitude, but see \chapref{\S 2.4}{02_testing_a_signal} \\ \hline
\end{tabularx}

The three differ in how much of the signal's magnitude they keep. They also differ in how
\emph{often} they trade, and that second axis is not a property of the rule alone: a crossover flips
whenever the two legs touch, and how often they touch is set by how large the market's moves happen
to be this month. Every parameter above is now chosen, and that count is still not under their
control.


\section{Volatility clustering breaks both}\label{sec:vol-clustering}

\begin{definition}[Volatility clustering]\label{def:vol-clustering}
Large moves are followed by large moves and small by small, \textbf{regardless of sign}. Returns
themselves are close to uncorrelated from one day to the next, but their \emph{absolute} values are
strongly and persistently autocorrelated. Direction is unpredictable; \textbf{amplitude is not}---a
market has stretches, days to months, where everything is simply bigger, without the underlying
trend having changed at all.
\end{definition}

\textbf{Neither piece survives it.} A short EWMA half-life and MACD's fast leg are the same object
under two names---a short average of recent returns---and a short average's swings scale with the
amplitude around it. Enter a cluster and the sign starts flipping repeatedly, not because the trend
turned but because the noise grew.

\fig{whipsaw-and-smoother}{One price path, calm at both ends with a shaded volatility cluster in the
middle, carrying a slow moving average and a fast one, every crossing marked. Left, the raw fast leg
crosses the slow leg nine times inside the cluster. Right, the same fast leg passed through a short
smoother first crosses three times. Illustrative.}{fig:whipsaw-and-smoother}

\textbf{The crossings are not merely expensive, they are backwards.} A downward spike drags the fast
leg under the slow one and flips the book short---at the bottom of the spike. The snap-back flips it
long again, at the top. Repeat that through a cluster and the rule has been buying high and selling
low on a schedule set by noise, which is worse than holding nothing: \textbf{whipsaw is not scatter
around the right answer, it is the opposite of the right answer}, taken over and over. Nine
crossings instead of three is nine round trips instead of three, and they land in exactly the
stretch where spreads are widest and depth is thinnest---at forty crossings a year and five basis
points a round trip, that is 200 bp of annual drag against a gross edge that might be four hundred.

\textbf{The local fix is a shorter second average, or a margin.} A \textbf{smoother}---a second
average applied to the signal after it is computed, shorter than the signal's own period or it
becomes another slow leg---stops a single day's noise from moving the sign on its own; MACD's 9-day
signal line is one (\secref{sec:three-series}). A \textbf{deadband}, which requires a crossing to
clear some margin before the position flips rather than acting on the exact touch, is the other.
Both read the signal and nothing else, so both are local: neither knows the tape has turned violent,
each damps every date alike, and both pay on the quiet days to protect the loud ones. Going further
means measuring the state of the market itself and asking what the rule is worth in each of its
values.

\pointer{\chapref{4 \midot{} Volatility Regimes}{04_volatility_regimes}.}


\section*{Appendix \midot{} Notation}
\addcontentsline{toc}{section}{Appendix \midot{} Notation}

Throughout, $t$ is the date. Everything in this chapter is computed one asset at a time, so the
asset subscript $s$ of \chapref{2}{02_testing_a_signal} is dropped.

\begin{xltabular}{\linewidth}{|>{\raggedright\arraybackslash}p{0.20\linewidth}|X|>{\raggedright\arraybackslash}p{0.16\linewidth}|}
\hline
\textbf{Symbol} & \textbf{Means} & \textbf{First used} \\ \hline
\endhead
$a_t$, $\lambda$
  & the newest observation, and the EWMA decay that keeps $\lambda$ of the old average
  & \secref{sec:ewma-recursion} \\ \hline
$H$
  & EWMA half-life, in periods
  & \secref{sec:half-life} \\ \hline
$N_f$, $N_s$
  & fast and slow lookback lengths, in periods
  & \secref{sec:pairing} \\ \hline
$P_t$, $\Delta_{t-j}$
  & price on that date, and the one-period change at that lag
  & \secref{sec:three-series} \\ \hline
$n_f$, $n_s$
  & fast and slow EMA spans, conventionally 12 and 26
  & \secref{sec:three-series} \\ \hline
$\alpha_f$, $\alpha_s$
  & their smoothing constants, two over span plus one
  & \secref{sec:three-series} \\ \hline
$c_i$, $k_j$
  & MACD's net weight on the price at that lag, and its kernel weight on the price change
  & \secref{sec:macd-is-momentum} \\ \hline
\end{xltabular}

\begin{note}[Collisions to watch]
$N_f$, $N_s$ are plain lookback lengths and $n_f$, $n_s$ are the EMA spans that play the same two
roles---case is the tell. $k_j$ is a kernel weight, not a portfolio weight:
\chapref{2}{02_testing_a_signal} reserves $w_{s,t}$ for the latter, which is why the kernel is not
written $w$. And $\alpha$ here is a smoothing constant, not
\chapref{\S 2.3.1}{02_testing_a_signal}'s \texttt{alpha} plotting keyword and not a regression
intercept; code font against maths is the tell.
\end{note}
